\section{Method}\label{sec:method}

\odrift{} follows the observability asymmetry of \cref{sec:setting}: decisions that need to
see small reward changes are made on history, and the anchor is used only where a coarse
decision suffices. It has four steps: segment the history by reward drift, fit joint
reward models, form one drift-corrected estimate per segment, and pool them with the anchor.

\subsection{A reward-drift monitor}
We fit one pooled per-action logistic model $\tilde q$ with linear and quadratic context
terms on all historical windows. For a block $B$ of windows and an action $a$, let
$m_B(a)$ be the mean residual $r_i-\tilde q(x_i,a)$ over records in $B$ with $a_i=a$, and
$v_B(a)$ its binomial variance estimate. Two adjacent blocks $L,R$ are compared by
\begin{equation}\label{eq:monitor}
  \chi^2_{LR}=\sum_{a}\frac{\big(m_L(a)-m_R(a)\big)^2}{v_L(a)+v_R(a)},
\end{equation}
referred to a $\chi^2$ distribution with one degree of freedom per action with enough
records. Binary segmentation splits the history at the most significant boundary while
the $p$-value is below $\alpha$, giving contiguous segments $S_1,\dots,S_J$ ordered from
newest to oldest.

\begin{proposition}[Invariance to context and logger drift]\label{prop:invariance}
Suppose $q_k=q$ for all windows in $L\cup R$ and $\tilde q\to q$. Then for every action
$a$, $\mathbb{E}[m_L(a)-m_R(a)]\to0$ whatever the context laws $P_k$ and loggers $\mu_k$,
so \eqref{eq:monitor} is asymptotically $\chi^2$ and the test has level $\alpha$ under
any context or logger drift. The mean-reward statistic $\bar r_L-\bar r_R$ and the
per-action raw-reward statistic do not have this property.
\end{proposition}
\begin{proof}[sketch]
In window $k$, $\mathbb{E}[R-q(X,a)\mid A=a,k]=\mathbb{E}_{P_k(x\mid a)}[q_k(x,a)-q(x,a)]=0$,
because $P_k(x\mid a)\propto P_k(x)\mu_k(a\mid x)$ only reweights a function that is
identically zero. Only $q_k\neq q$ moves $m_B(a)$. In contrast,
$\mathbb{E}[R\mid k]=\sum_a\int P_k(x)\mu_k(a\mid x)q(x,a)\,dx$ changes with $P_k$ and
$\mu_k$ when $q$ is fixed, which gives false alarms; and a reward change can be cancelled
by an opposite change of context or logger, which gives missed drift. The per-action raw
mean $\mathbb{E}_{P_k(x\mid a)}[q(x,a)]$ removes the action mix but still depends on
$P_k$ and on $\mu_k$ as a function of $x$.
\end{proof}
If $q$ lies outside the model class, the residual has conditional mean
$\mathbb{E}_{P_k(x\mid a)}[q-\tilde q]$, which depends on $P_k$; context drift can then
trigger the monitor. The richer the model class, the smaller this effect.

\subsection{Joint segment reward models}
Reward models for estimation are cross-fitted~\cite{chernozhukov2018double} per-action
logistic regressions on a context basis $\phi(x)$. Across segments we fit one joint model
\begin{equation}
  \operatorname{logit} \hat q_j(x,a)=\phi(x)^\top\beta_a+\phi(x)^\top\delta_{j,a},\qquad \delta_{1,a}=0,
\end{equation}
a shared base plus per-segment offsets, with an $L_2$ penalty on the offsets scaled by
$1/s^2$. The scale $s\le1$ acts as a ridge prior on drift size and interpolates between
fully corrected ($s=1$) and uncorrected reuse ($s\to0$). The basis is chosen per episode
between linear features and $D=32$ random Fourier features~\cite{rahimi2007random} by
held-out log-loss of the pooled model; the same basis is used by every method.

\subsection{Reference segment and drift-corrected reuse}
The reference segment is the newest segment whose uncorrected estimate the anchor does
not contradict at $|z|\le z_{\mathrm{anc}}$; if all are contradicted the anchor is used
alone. Each segment $S_j$ contributes
\begin{equation}\label{eq:segment}
  \hat V_j=\frac1M\sum_{m=1}^M\sum_a\pi(a\mid x^*_m)\hat q_{\mathrm{ref}}(x^*_m,a)
  +\frac1{n_j}\sum_{i\in S_j}W_i\,\big(r_i-\hat q_j(x_i,a_i)\big),
\end{equation}
\begin{equation}
  W_i=\frac{\pi(a_i\mid x_i)}{\sum_{k\in S_j}\frac{n_k}{n_j}\,\mu_k(a_i\mid x_i)/w_k(x_i)}.
\end{equation}
The direct-method term uses the reference model, so it targets $q_{\mathrm{ref}}$; the
correction term uses the segment's own model, so it removes the segment's own model error.
$W_i$ is the balanced importance weight from the target to the segment's mixture of
loggers and context laws~\cite{agarwal2017effective,kallus2021optimal}, with $w_k$
estimated by a probabilistic classifier. For the reference segment $\hat q_j=\hat q_{\mathrm{ref}}$
and \eqref{eq:segment} is the usual balanced DR estimate. For an older segment, the
difference $\hat q_{\mathrm{ref}}-\hat q_j$ is exactly the fitted reward drift, so
\eqref{eq:segment} transfers the segment's evidence to the reference reward when (A4) holds.

\subsection{Pooling and interval}
The anchor DR estimate joins the segment estimates, and the weights
$\hat\alpha=\arg\min_{\alpha\ge0,\,\mathbf1^\top\alpha=1}\alpha^\top\hat\Sigma\alpha$ minimise
the estimated variance, where $\hat\Sigma$ accounts for the target contexts shared by
all direct-method terms. Because bias may remain, the reported interval is
DerSimonian--Laird~\cite{dersimonian1986meta}, inflating the standard error by the
between-component heterogeneity $\hat\tau^2$; a sample-split interval with an explicit bias
allowance is reported as a certified but wide alternative.
Without an anchor \odrift{} returns an estimate flagged as not certified.


\paragraph{Settings.} $\alpha=10^{-4}$, $z_{\mathrm{anc}}=5$ and $s=0.25$ were tuned on development data (\cref{sec:setup}).
